% !TEX root = ../main.tex
\clearpage
\phantomsection
\section*{1 \textbf{서론 및 배경}}
\label{sec:introduction-and-background}
\addcontentsline{toc}{section}{1 서론 및 배경}

탈중앙화 금융(DeFi)은 공개 블록체인상에서 대출, 차입, 레버리지 거래 활동을 확대해 왔으나, 대부분의 신용 유사 상호작용은 여전히 높은 담보 비율 하에서 이루어진다. 이는 프로토콜이 전통적인 오프체인 신원, 고용, 신용 이력에 의존할 수 없기 때문이다(Sch\"{a}r, 2021; Packin \& Lev-Aretz, 2024). 가명(pseudonymous) 환경에서 핵심 과제는 단순히 자본을 배분하는 것이 아니라 안전하게 배분하는 것이다: 신용 유사 상호작용(예: 반복적 차입, 상환, 지출 권한 위임, 레버리지 노출의 성공적 관리)에서 신뢰성 있게 행동할 가능성이 높은 지갑을 식별하고, 기회주의적·악의적 행위자(예: 포지션 청산, 그리핑, Sybil식 계정 생성, 급속한 평판 농사)에 대한 노출을 제한하는 것이다(Packin \& Lev-Aretz, 2024).


본 연구에서 \textbf{온체인 평판 점수(on-chain reputation score)} 산출이란, 관찰 가능한 온체인 행동과 관계를 사용하여 블록체인 주소에 대한 정량적 점수를 계산하는 과정을 말하며, 오프체인 개인 데이터에 의존하지 않고 신뢰성 또는 신용 관련 신뢰도를 근사하는 것을 목표로 한다(Ghosh et al., 2024; Vathsala V. \& V. D. Sharma, 2025). DeFi에서 이러한 점수는 (i) 온체인 신용 리스크(예: 포지션 청산 또는 반복적 거래 손실)의 보다 정확한 평가, (ii) 대출 및 무기한 거래 프로토콜의 리스크 통제 개선, (iii) 저위험 참여자에 대한 담보 요건 완화 등 자본 효율적 신용으로의 잠재적 경로를 지원할 수 있다(Packin \& Lev-Aretz, 2024).


모호함을 피하기 위해, 본 계획서 전반에서 사용하는 핵심 개념을 다음과 같이 정의한다.


\begin{itemize}
\item[\textbullet] \textbf{DeFi(Decentralized Finance): }공개 블록체인의 스마트 계약을 통해 구현되는 금융 서비스(예: 대출, 차입, 무기한 거래)(Sch\"{a}r, 2021).
\item[\textbullet] \textbf{스마트 계약(smart contract): }블록체인에 배포된 자동 실행 프로그램 코드로, 사전 정의된 조건이 충족되면 신뢰할 수 있는 중개자 없이 계약 조건을 자동으로 집행한다(Buterin, 2014; Szabo, 1997).
\item[\textbullet] \textbf{지갑/주소(Wallet / Address): }자산을 보유하고 스마트 계약과 상호작용할 수 있는 블록체인 계정(분석 단위로 사용)(Buterin, 2014; Ethereum Improvement Proposals, 2015).
\item[\textbullet] \textbf{온체인 평판 점수: }금융 상호작용에서 지갑의 신뢰성을 나타내기 위해 온체인 신호로부터 계산된 수치 지표(Do et al., 2023).
\item[\textbullet] \textbf{온체인 신용 리스크: }공개 블록체인에서 관찰 가능한 부정적 결과, 예를 들어 포지션 청산, 반복적 거래 손실, 프로토콜이 정의한 연체 신호(Lin et al., 2024).
\item[\textbullet] \textbf{GMX V2 무기한 스왑 포지션: }Arbitrum One의 GMX V2 \textbf{EventEmitter}가 방출하는 \textbf{PositionIncrease} 및 \textbf{PositionDecrease} 이벤트를 통해 개설·청산되는 GMX V2의 담보화된 롱 또는 숏 노출(GMX, n.d.; Arbitrum Foundation, n.d.).
\item[\textbullet] \textbf{수익성 청산/청산 성공(profitable close / close-success): }포지션 청산이 아닌(\texttt{orderType} $\neq 7$) 양의 실현 손익(\textbf{basePnlUsd} $> 0$)을 가진 \textbf{PositionDecrease}(GMX, n.d.).
\item[\textbullet] \textbf{청산 성공 횟수(close-success count): }관찰 기간 내 지갑의 수익성 \textbf{PositionDecrease} 이벤트 수(송금 기반 PageRank 검증의 in-degree에 유사)(GMX, n.d.; Do et al., 2023).
\item[\textbullet] \textbf{실현 이익 대리변수(realized gain proxy): }관찰 기간 내 수익성 청산에 대한 $\max(\text{basePnlUsd}, 0)$의 합(in-value에 유사)(GMX, n.d.; Do et al., 2023).
\item[\textbullet] \textbf{가중 PageRank(Weighted PageRank): }엣지 강도로 가중된 방향성 엣지를 통해 영향력을 전파하는 그래프 기반 순위 방법; 본 연구는 전 과정에서 가중 변형을 사용한다(Page et al., 1998; Do et al., 2023).
\item[\textbullet] \textbf{ERC-20 승인(allowance): }토큰 소유자가 지출자(spender)가 대신 이체할 수 있는 금액을 기록하는 상태 변수(Ethereum Improvement Proposals, 2015).
\item[\textbullet] \textbf{Approve / Permit: }승인을 설정하는 ERC-20 메커니즘; permit(EIP-2612)은 서명 메시지를 통해 승인을 갱신할 수 있게 한다(Ethereum Improvement Proposals, 2020).
\item[\textbullet] \textbf{보증 그래프(endorsement graph): }단순 송금이 아닌 명시적 위임 또는 보증 신호를 나타내는 엣지로 구성된 방향 그래프(Do et al., 2023).
\end{itemize}
신용 기반 및 레버리지 DeFi에서 평판 점수 산출은 신뢰할 수 있는 참여자와 잠재적 악의적 행위자를 구별하는 데 중요하다(Sch\"{a}r, 2021; Packin \& Lev-Aretz, 2024). 온체인 평판 점수 산출에 대한 최근 PageRank 영감 접근법에는 Do et al.(2023)이 제안한 AWP 알고리즘이 포함되며, 이는 거래량(유입이 클수록 신뢰가 높음)과 활동성(최근 거래에 로지스틱 시간 감쇠 함수로 더 높은 가중치 부여)에 기반하여 평판 점수를 부여한다. 이 접근은 경제적 규모와 시간적 관련성을 모두 포착하지만, 명시적 지출 승인을 신뢰 신호로 직접 모델링하지는 않는다.


그러나 ERC-20 토큰 표준에는 토큰 보유자가 다른 주소에 대해 지정된 금액까지 대신 지출할 수 있는 권한을 부여하는 두 가지 명시적 승인 메커니즘---approve와 permit---이 포함되어 있다(Ethereum Improvement Proposals, 2015; Ethereum Improvement Proposals, 2020). 각 승인 이벤트는 본질적으로 보증을 나타내며, 승인된 금액이 신뢰의 세분화된 정도를 직접 정량화한다. 각 소유자\(\to\)지출자 쌍에 대해 최신 승인만 추출하면, 각 엣지의 가중치가 가장 최근 승인된 토큰 금액과 같은 방향성 가중 보증 그래프 $G_{\text{approve}}$를 구축할 수 있다(Do et al., 2023). 이 설계는 새 승인이 이전 승인을 덮어쓰므로 시간 기반 감쇠가 필요 없으며, 가장 현저하고 최신의 신뢰 신호에 계산을 집중한다.


\textbf{GMX V2에서의 도메인 검증.} 본 연구의 주요 실증 설정은 \textbf{Arbitrum One}과 \textbf{GMX V2 무기한 스왑 포지션 거래}이다. 원칙적으로 담보 부족 대출의 부도 또는 청산 이벤트가 가장 직접적인 신용 리스크 레이블을 제공하지만, Arbitrum One에서는 본 학위논문에 필요한 규모의 안정적 평가 데이터셋을 구성하기에 관찰 가능하고 재현 가능한 무담보 또는 담보 부족 대출 활동의 양이 불충분하다. 따라서 동일 체인에서 밀도 높은 지갑 수준 실현 결과를 제공하는 GMX V2 거래 성공 대리변수---\textbf{PositionDecrease}를 통한 롱·숏 무기한 포지션 청산---에 대해 EndorseRank를 검증한다. 무기한 포지션 개설은 담보화된 레버리지 노출이다: 트레이더는 마진을 예치하고 기초 자산에 대한 방향성 리스크를 부담한다. 양의 \textbf{basePnlUsd}---수익성 \textbf{PositionDecrease}---로 포지션을 청산하는 것은 신용 맥락에서 차입과 상환을 성공적으로 관리하는 것과 기능적으로 유사하다. AWP 논문의 in-degree 및 in-value를 도메인 직관 대리변수로 사용한 것을 따르면, 더 높은 평판 점수가 더 빈번한 수익성 청산(\textbf{청산 성공 횟수})과 더 큰 누적 실현 이익(\textbf{실현 이익 대리변수})을 달성하는 지갑과 정렬될 것으로 기대한다.


DeFi에서의 중요성을 명확히 하면, 온체인 평판 점수 산출의 향상은 다음과 같은 구체적 이점을 제공할 수 있다.


\begin{itemize}
\item[\textbullet] \textbf{리스크 관리 및 프로토콜 안전: }위험한 지갑의 더 나은 식별은 부실채권을 줄이고 대출 및 무기한 거래 시장의 안정성을 개선할 수 있다(Lin et al., 2024).
\item[\textbullet] \textbf{자본 효율: }더 유익한 신뢰 신호는 일률적 제약 대신 차별화된 리스크 파라미터(예: 차입 한도 또는 담보 요건)를 지원할 수 있다(Packin \& Lev-Aretz, 2024).
\item[\textbullet] \textbf{오프체인 데이터 없는 접근 및 포용: }온체인 점수 산출은 전통적 신용 기록이 없는 사용자의 참여를 가능하게 하면서 DeFi의 무허가 접근 원칙을 유지할 수 있다(Sch\"{a}r, 2021).
\item[\textbullet] \textbf{DeFi 응용을 위한 조합 가능한 신뢰: }재사용 가능한 평판 계층은 여러 프로토콜(예: 대출자, 무기한 거래소, 보험, DAO)이 일관되고 설명 가능한 결정을 내리는 데 활용될 수 있다(Packin \& Lev-Aretz, 2024).
\end{itemize}
명확한 보증 의미론과 승인 기반 관계의 더 희소한 그래프 구조에도 불구하고, ERC-20 승인을 PageRank의 주요 엣지 가중치로 체계적으로 활용한 기존 연구는 없다. 이 방법론적 공백이 $G_{\text{approve}}$에 표준 Weighted PageRank를 적용하는 본 연구의 제안 온체인 평판 알고리즘 EndorseRank의 동기가 된다. EndorseRank는 GMX 거래 성공 대리변수와 동등하거나 우수한 정렬을 달성하면서 거래 가치 및 시간 감쇠 기반 방법 대비 데이터 처리 복잡도를 크게 줄일 것으로 가설 설정한다. 이 동기와 연구가 제공할 결과 간의 명확한 정렬을 위해, 본 계획서는 목표 연계 기대 결과를 명시한다.


\begin{itemize}
\item[\textbullet] Arbitrum One에서 ERC-20 approve/permit 상태로부터 파생된 재현 가능한 파이프라인 및 최신 승인 보증 그래프;
\item[\textbullet] 이 희소하고 신뢰 의미론적 그래프에서 동작하는 최적화되고 확장 가능한 Weighted PageRank 구현(EndorseRank); 및
\item[\textbullet] GMX V2 \textbf{PositionDecrease} 결과와 순위 상관 지표(Spearman's rho, Kendall's tau)를 사용하여 EndorseRank가 거래 기반 기준선 대비 거래 성공 정렬에서 동등하거나 개선된 정렬을 보이면서 계산 부담을 줄인다는 실증적 증거.
\end{itemize}
